% RNAAS research note, AASTeX. Draft written by Claude (AI) for the researcher to review
% and edit. Results are placeholders until the pre-registered analysis has been run on
% Gaia DR4 (released 2 Dec 2026). Word limit: 1500 words incl. title, captions, references.
% NOTE: the class/option below must be checked against the current AASTeX release
% (AASTeX v7 is preferred by the AAS; the RNAAS option is documented for v6.3.1).
\documentclass[RNAAS]{aastex631}

\usepackage{amsmath}

% ---- numbers filled from the repository (update before submission) ----
\newcommand{\Ncand}{67}            % candidates in the scored samples
\newcommand{\Ntest}{[N\_TESTABLE]}        % testable candidates
\newcommand{\NG}{111}                     % Graham+2015
\newcommand{\NC}{33}                      % Charisi+2016 final sample
\newcommand{\NF}{257}                     % Foustoul+2025 (36 + 221)
\newcommand{\prereghash}{[PREREG\_COMMIT\_HASH]}
\newcommand{\codehash}{9284e83}
\newcommand{\pending}[1]{\textbf{[PENDING: #1]}}

\begin{document}

\title{A Pre-registered Test of Published Periodic-Quasar Ephemerides with Gaia DR4 Epoch Photometry}

\author{[AUTHORS --- TBD]}

\begin{abstract}
Several hundred quasars have been proposed as sub-parsec supermassive black hole binaries
because their optical light curves appear periodic over 1.5--3 cycles. Red-noise
variability produces such apparent periods readily. We test whether the published periods
of the CRTS \citep{Graham2015} and PTF \citep{Charisi2016} candidates, with phase and
amplitude taken from the discovery data alone, predict Gaia DR4 G-band photometry
obtained after each catalogue was published better than a damped random walk conditioned
on the same data. The models, frozen parameters, scoring rule and decision thresholds were
committed to a public repository before the DR4 release (commit \prereghash).
\pending{headline result: H1 outcome, number of survivors, 95\% upper limit on the
fraction of true periodic sources.}
\end{abstract}

\section{Introduction}

Supermassive black hole binaries (SMBHBs) are expected after galaxy mergers and are the
likely source of the nanohertz gravitational-wave background reported by pulsar timing
arrays \citep{NANOGrav2023}. At sub-parsec separations their orbital periods are years,
and accretion modulated by the binary, or relativistic Doppler boosting, could make the
quasar light curve periodic. Searches of large time-domain surveys have produced hundreds
of candidates \citep{Graham2015, Charisi2016, Liu2019, Chen2024, Foustoul2025}. Most show
only 1.5--3 cycles, and stochastic quasar variability, well described by a damped random
walk \citep[DRW;][]{Kelly2009}, produces apparent periods of this kind at a high rate
\citep{Vaughan2016}. Extended baselines have removed most candidates
\citep{Liu2016, Liu2019, Foustoul2025}, Gaia DR3 sinusoids failed to predict later ZTF
data \citep{ElBadry2026}, and a blind search of Gaia DR3 AGN found no reliable periodic
source \citep{Huijse2026}.

The cleanest test of a claimed period is whether it predicts data that played no part in
finding it. Gaia DR4 publishes epoch photometry for all sources from 2014 July 25 to 2020
January 20, which fills the gap between CRTS (to $\sim$2013) and ZTF (from 2018). For the
CRTS and PTF candidates, the DR4 epochs after each paper's submission are genuinely new.
We fixed every prediction and decision rule before DR4 was released, so that the analysis
could not be tuned to the outcome.

\section{Pre-registration and data}

The pre-registration (repository commit \prereghash, made before 2026 December 2) contains
the candidate list with the source of every value, the analysis code (commit \codehash),
the frozen model parameters and predictions for every candidate, the scoring rule and the
decision thresholds. Any later deviation is reported as a dated amendment.

\textit{Candidates.} The primary (forecast) sample comprises the \NG\ CRTS candidates of
\citet{Graham2015} and the \NC\ PTF candidates of \citet{Charisi2016}. A secondary sample,
the 36 strong and 221 weaker candidates of \citet{Foustoul2025}, is tested in the gap
between their CRTS and ZTF data. The PS1 candidates of \citet{Liu2019} (discovery light
curves not public) and the three ZTF candidates of \citet{Chen2024} (extended-baseline
periods not tabulated) are listed but not scored.

\textit{What is frozen.} No source paper published a phase or reference epoch, and
\citet{Graham2015} and \citet{Charisi2016} gave no period uncertainties. We therefore
freeze each published period exactly as printed, and derive the phase and amplitude only
from discovery data obtained before the paper's arXiv v1 submission date (the GP posterior
described below). These values were frozen before any DR4 data existed. The test is thus a
test of the published period together with the phase implied by the discovery data.

\textit{Data.} Discovery light curves are CRTS (Catalina Sky Survey photometry from the
Catalina Surveys data release; unfiltered, V-calibrated) for the Graham sample and PTF/iPTF
Mould-$R$ \citep[IRSA;][]{Law2009} for the Charisi sample, binned by night. The CRTS service was
intermittent in November 2026; Graham candidates whose CRTS light curve could not be retrieved
before the freeze are conditioned on PTF/iPTF data before the cutoff instead and reported
separately (41 objects; 36 more have neither). For each candidate, the Gaia test epochs are the DR4 G-band
field-of-view transits after the arXiv v1 date of the source paper (2015 July 27 and 2016
April 4), binned into visits of $\leq$1\,d. Candidates must have Gaia DR3 mean $G<20.7$.
Nobody on the project inspected Gaia DR3 epoch photometry, which covers the first 34
months of the DR4 window for some AGN; a robustness test uses only DR4 epochs after 2017
May 28.

\section{Method}

Both models describe the discovery-band magnitude as a Gaussian process. Model D is a DRW
with covariance $\sigma^2\exp(-|\Delta t|/\tau)$ plus a mean level. Model S adds
$a\cos(2\pi t/P)+b\sin(2\pi t/P)$ at the published period $P$, with a broad Gaussian prior
($\sigma_s=1$\,mag) on $a,b$, so that the published sinusoid is taken as given and its phase
and amplitude are set by the discovery data. Hyperparameters ($\sigma,\tau$ and an extra
white-noise term) are fitted to the discovery data only and frozen. Conditioning on the
discovery data gives each model's predictive distribution at the Gaia epochs; in model S
this carries the posterior phase and amplitude of the sinusoid. Gaia magnitudes are
modelled as $c+s\,x(t)+\epsilon$. The zero-point offset $c\sim\mathcal{N}(0,1\,\mathrm{mag}^2)$,
the amplitude ratio $s\sim\ln\mathcal{N}(0,0.35^2)$ and the extra Gaia noise
$j\sim$ half-normal(0.05\,mag) are marginalised identically in both models. The score is the
predictive log Bayes factor
\begin{equation}
\ln B=\ln p(\mathbf{y}_{\rm G}\mid D_{\rm disc},S)-\ln p(\mathbf{y}_{\rm G}\mid D_{\rm disc},D).
\end{equation}

Because both models condition on the actual discovery data, the null distribution of
$\ln B$ for each candidate is obtained by simulating Gaia data from model D's predictive
distribution at the real DR4 epochs (1000 draws); the candidate's selection, which depended
only on the discovery data, is thereby accounted for. Simulating from model S gives the power.
A candidate is testable if, using the DR4 timestamps and uncertainties only, at least half of
the model-S simulations reach $p<0.05$.

The plug-in hyperparameters make these $p$-values slightly optimistic. Before release, we
simulated DRW ``fake candidates'' on each forecast candidate's cadence. Each fake was required
to look at least as periodic at the published period as the real candidate, had all
hyperparameters refitted, and was scored exactly as the data. This gave a per-candidate
false-positive rate $r_0=0.056$ (13,400 fakes; 95\% upper bound 0.060) instead of 0.05. A steeper
(damped harmonic oscillator) red-noise process gave $r_0=0.051$ for the PTF sample.

Decision rules (frozen): (H1) the population shows an excess if the number $k$ of testable
forecast candidates with $p<0.05$ has a one-sided binomial probability $<0.01$ under rate
$\max(0.05,r_0^{95})$, where $r_0^{95}$ is the 95\% upper bound on $r_0$; (H2) survivors have
$\ln B>3$ and Benjamini--Hochberg $q<0.05$, and are called ephemeris-specific if their $\ln B$
also exceeds that of 95\% of 200 phase-shifted forecasts (H4); (H3) the fraction $f$ of true
periodic sources is constrained by a mixture of each candidate's null and injection $\ln B$
distributions, and we report its 95\% upper limit. A dress rehearsal on mock DR4 data recovered
$H_0$ when the data were drawn from model D and rejected it when they were drawn from model S.

\section{Results}

\pending{Number of candidates with DR4 counterparts and enough test epochs; number
testable.}

\pending{Figure 1: distribution of $\ln B$ for testable candidates compared with the stacked
null and injection distributions (left); posterior on $f$ (right).}

\pending{H1 count $k$ vs. expected; survivors with $q$ values and H4 outcomes; $f$ upper
limit; sensitivity analyses (steeper-spectrum null, DR4-only epochs, unbinned data, period
marginalisation, wider amplitude-ratio prior).}

\section{Discussion}

\pending{Comparison with \citet{Foustoul2025} and \citet{ElBadry2026}; implications for the
binary population and the PTA background \citep{Sesana2018}; limitations: sinusoid-only
model \citep{Lin2026}, sparse Gaia cadence, filter differences.}

\begin{acknowledgments}
This work uses data from the European Space Agency mission Gaia, processed by the Gaia Data
Processing and Analysis Consortium (DPAC); from the Catalina Real-time Transient Survey;
and from the Palomar Transient Factory and the Zwicky Transient Facility, obtained from the
NASA/IPAC Infrared Science Archive. [AI-USE STATEMENT --- TBD]
\end{acknowledgments}

\software{astropy, astroquery, celerite2 \citep{celerite2}, numpy, scipy, pandas}

\bibliographystyle{aasjournal}
\bibliography{refs}

\end{document}
